\section{Background}\label{sec:background-and-context}

Pests and plant diseases cause substantial crop losses, so crop protection is important for agricultural production, so it's necessary to eradicate both plagues and plant diseases. 
Many products use chemical substances, which are harmful to both the environment and the very crop they are trying to treat, and as such we need to bet on biological control agents (BCAs).
BCAs can contribute to integrated pest management as alternatives. 
Their practical use depends on the crop, target organism, microbial strain, formulation, application conditions, and regulatory context. 
Microbial agents can act through several mechanisms, including plant defence induction, competition, hyperparasitism, and antibiosis. 
These mechanisms and their dependence on the local environment make evidence from an authorisation record different from evidence of efficacy in a new field setting.

Recommender Systems (RS) are software tools used for recommending the most suitable items based on a user's profile. 
RS have been largely applied to fields such as movies, music, and books.
Recently, they have been studied as alternatives to predict links in scientific items, such as chemical compounds, genes and BCAs. 
The project PlantPathoRec \cite{plantpathorec-project}, whose goal is to create a platform for both researchers and farmers, intends to recommend new and commercial BCAs to treat plant diseases, using RS approaches. 
This projects goal is to contribute to the recommendation of commercial BCAs.

\subsection{Related Work}

Previous work describes the opportunities and adoption barriers for microbial and other biological control agents, including variable field performance, formulation constraints, and the need to integrate them into broader pest-management systems. 
Research on microbial modes of action also shows that disease suppression can result from interacting mechanisms rather than one universal agent--pathogen effect. 
These findings motivate a system that makes existing authorisation associations easier to retrieve and inspect while avoiding claims about biological performance. 
The present work, therefore, focuses on data integration and ranking of traceable candidates.